\documentclass{article}

\usepackage{amsmath,amssymb}
\usepackage{xurl}
\usepackage[colorlinks=true,linkcolor=blue,citecolor=blue,urlcolor=blue]{hyperref}
\setlength{\emergencystretch}{2em}

% Shared commands for notes in docs/paper-gaps/.
% Notation follows the LDT paper (references/ldt-paper/).

\newcommand{\Lean}{\textsc{Lean}}
\newcommand{\leanid}[1]{\textsf{#1}}
\newcommand{\ghissue}[1]{\href{https://github.com/LionSR/MIPStarRE/issues/#1}{GitHub issue \##1}}
\newcommand{\ghpr}[1]{\href{https://github.com/LionSR/MIPStarRE/pull/#1}{GitHub PR \##1}}

% --- Paper-compatible notation ---
\newcommand{\F}{\mathbb{F}}
\newcommand{\N}{\mathbb{N}}
\newcommand{\C}{\mathbb{C}}
\newcommand{\tr}{\operatorname{tr}}
\newcommand{\ot}{\otimes}
\newcommand{\E}{\mathop{\bf E\/}}
\newcommand{\bx}{\mathbf{x}}
\newcommand{\by}{\mathbf{y}}
\newcommand{\bu}{\mathbf{u}}
\newcommand{\bv}{\mathbf{v}}

% Approximation relation (paper uses \approx_\eps and \simeq_\zeta)
\newcommand{\approxeps}{\approx_{\varepsilon}}
\newcommand{\simgeq}{\simeq_{\zeta}}

% Polynomial spaces (paper uses \polyfunc, \polysub, \polymeas)
\newcommand{\polyfunc}[3]{\operatorname{Poly}(#1,#2,#3)}
\newcommand{\polysub}[3]{\operatorname{PolySub}(#1,#2,#3)}
\newcommand{\polymeas}[3]{\operatorname{PolyMeas}(#1,#2,#3)}


\title{The ambient dimension in the commutativity theorem for $G$}
\author{}
\date{2026-05-01}

\begin{document}
\maketitle

\section{The assertion}

We examine the statement of the commutativity theorem for the slice
measurements $G^x$ in \cite[Section~\emph{Commutativity of~$G$}]{Ji2020LowIndividualDegree}.%
\footnote{Tracked in \ghissue{930}. The paper passage is
\path{references/ldt-paper/commutativity-G.tex}, lines~228--257.}
The source states that $(\psi,A,B,L)$ is an
$(\eps,\delta,\gamma)$-good symmetric strategy for the $(m,q,d)$ low
individual degree test. It then assumes, on average over
$(\bu,\bx) \in \F_q^{m+1}$,
\[
  A^{u,x}_a \ot I \simeq_\zeta I \ot G^x_{[g(u)=a]},
\]
where $G^x$ is a projective submeasurement with outcomes in
$\polyfunc{m}{q}{d}$.

\section{The point at issue}

The displayed consistency condition is naturally a hypothesis about a
strategy for the $(m+1,q,d)$ test. In that test, the point measurement
$A$ is indexed by points of $\F_q^{m+1}$, which are written as
$(u,x)$ with $u\in\F_q^m$ and $x\in\F_q$. A strategy for the
$(m,q,d)$ test has point measurements indexed by $\F_q^m$ instead, so
$A^{u,x}_a$ is not the corresponding point measurement.

The surrounding argument also uses the theorem in the successor step:
the slice measurements $G^x$ live in dimension $m$, and they are pasted
to produce a measurement in dimension $m+1$. Thus the intended ambient
strategy is the $(m+1,q,d)$ strategy, while the slice family remains a
family of $m$-variable polynomial submeasurements.

\section{The formal and blueprint statements}

The blueprint states the corrected ambient dimension: in
\path{blueprint/src/chapter/ch08_commutativity.tex}, the theorem
\leanid{thm:com-main} assumes an $(\eps,\delta,\gamma)$-good symmetric
strategy for the $(m+1,q,d)$ low individual degree test, and takes
projective submeasurements $G^x$ in $\polysub{m}{q}{d}$.

The formal theorem uses the same convention.%
\footnote{The declaration is \leanid{MIPStarRE.LDT.Commutativity.comMain}
in \doclink{MIPStarRE/LDT/Commutativity/Main/Results.lean}. Its parameter
\leanid{params} describes the slice dimension $m$, while the strategy has
type \leanid{SymStrat params.next}, where \leanid{params.next.m = params.m+1}.}
It proves the commutativity conclusion for the full slice products
\[
  G^x_g G^y_h \ot I \approx_\nu G^y_h G^x_g \ot I,
  \qquad
  \nu = 30m\left(\gamma^{1/4}+\zeta^{1/4}+(d/q)^{1/4}\right),
\]
with the same three hypotheses of consistency with $A$, strong
self-consistency, and boundedness.

\section{Conclusion}

The occurrence of $(m,q,d)$ in the source statement of
\leanid{thm:com-main} is a harmless dimension typo. The theorem is used
with an ambient $(m+1,q,d)$ strategy and a family of $m$-dimensional slice
measurements. The blueprint and the formal development use this corrected
statement. No Lean theorem statement or proof needs to change.

\bibliographystyle{alpha}
\bibliography{references}

\end{document}